\chapter{Itchy Skin (Pruritus)}
\index{Itchy Skin (Pruritus)}

\section*{Definition and Overview}
Itchy skin, medically called pruritus, is an unpleasant sensation that triggers the urge to scratch. It can be caused by skin conditions such as eczema, hives, and dry skin; by insect bites and infestations; by allergic reactions; and by internal conditions including liver, kidney, thyroid, and blood disorders, which can cause itching without a rash. Most itching settles with moisturising, cool compresses, and antihistamines, and treating the cause. Persistent itching that does not respond to simple measures, especially without a visible rash, should be assessed, because it may signal an underlying medical condition.

\section*{Epidemiology}
Itching is one of the most common symptoms in medicine, affecting most people at some time and around one in ten people chronically. Dry skin is the most common cause, particularly in older people, in winter, and with frequent hot showers. Skin conditions including eczema, which affects around one in ten people, and hives are common causes. Internal causes are less common but important: itching affects many people with chronic kidney disease, liver disease with cholestasis, iron deficiency, thyroid disease, and some cancers and blood disorders. Pregnancy-related itching is also recognised.

\section*{Aetiology and Pathophysiology}
Itch is transmitted by specialised nerve fibres in the skin to the spinal cord and brain, and it is amplified by histamine and other chemicals released by allergic and inflammatory cells. In dry skin, the damaged skin barrier exposes nerve endings and triggers itching. Eczema and hives involve inflammation and histamine release, while infestations such as scabies cause itch through the mites and their products. In systemic disease, itch arises from circulating substances: bile salts in liver disease, histamine and other mediators in some blood disorders, and metabolic waste products in kidney failure. The scratch-itch cycle, in which scratching damages the skin and releases more itch mediators, maintains and worsens itching.

\section*{Clinical Features}
\begin{itemize}
    \item Itching, which may be generalised or localised
    \item A visible rash, redness, dryness, or weals in skin-related causes
    \item Itching without a rash, suggesting a possible internal cause
    \item Scratch marks, excoriations, and broken skin
    \item Thickened, leathery skin from chronic scratching
    \item Dryness and flaking, especially in winter or with hot showers
    \item Itching that is worse at night or in bed
    \item Associated features pointing to the cause: jaundice, fatigue, weight loss, or fever
    \item Disturbed sleep and irritability
\end{itemize}

\section*{Differential Diagnosis}
The cause of itching is identified through the pattern of symptoms and investigations.
\begin{itemize}
    \item \textbf{Skin conditions:} dry skin, eczema, hives, psoriasis, and dermatitis
    \item \textbf{Bites and infestations:} insect bites, scabies, and lice
    \item \textbf{Allergic reactions:} to foods, medicines, or contact substances
    \item \textbf{Liver disease:} cholestasis with jaundice and dark urine
    \item \textbf{Kidney failure:} uraemic pruritus in advanced disease
    \item \textbf{Blood disorders:} iron deficiency, polycythaemia, and some cancers
    \item \textbf{Endocrine conditions:} thyroid disease and diabetes
    \item \textbf{Medicines:} drug reactions and opioid-related itch
    \item \textbf{Pregnancy:} specific conditions such as obstetric cholestasis
\end{itemize}

\section*{When to Seek Medical Care}
Medical assessment is appropriate for itching that is severe, persistent, or affecting sleep and quality of life, for itching without a rash, and for itching with weight loss, fatigue, fever, jaundice, or other systemic symptoms. Itching in pregnancy should always be reviewed, as it may indicate obstetric cholestasis. Seek urgent care for itching with breathing difficulty or facial swelling, which may indicate a severe allergic reaction.

\textbf{Contact your local emergency services for:}
\begin{itemize}
    \item Itching with swelling of the face, lips, or throat, or difficulty breathing, which may indicate anaphylaxis
    \item A widespread rash with fever and feeling unwell
    \item Rapidly spreading redness, pain, and swelling, which may indicate infection
    \item Collapse, confusion, or faintness
\end{itemize}

\section*{Diagnosis and Investigations}
\begin{itemize}
    \item \textbf{History:} onset, pattern, triggers, medicines, and associated symptoms
    \item \textbf{Examination:} the skin, nails, and features of underlying conditions
    \item \textbf{Blood tests:} full blood count, iron studies, thyroid, liver, and kidney function
    \item \textbf{Allergy assessment:} identification of triggers where allergy is suspected
    \item \textbf{Skin scrapings:} for scabies and fungal infection
    \item \textbf{Liver and bile duct imaging:} where liver disease is suspected
    \item \textbf{Specialist referral:} dermatology, and other specialists as indicated by the cause
\end{itemize}

\section*{Management}
\begin{itemize}
    \item \textbf{Moisturising:} regular emollients for dry skin, applied especially after bathing
    \item \textbf{Cool compresses and cool bathing:} to soothe the skin
    \item \textbf{Antihistamines:} for allergy-related and histamine-driven itching
    \item \textbf{Topical steroids:} for inflammatory conditions such as eczema
    \item \textbf{Avoiding triggers:} irritants, allergens, heat, and rough fabrics
    \item \textbf{Treatment of the cause:} iron replacement, treatment of liver or kidney disease, and thyroid therapy
    \item \textbf{Targeted treatments:} for scabies, fungal infection, and other specific causes
    \item \textbf{Specialist treatments:} phototherapy and systemic medicines for severe or refractory itching
    \item \textbf{Breaking the scratch-itch cycle:} short nails, distraction, and covering itchy areas
\end{itemize}

\section*{Complications}
\begin{itemize}
    \item Skin damage, excoriation, and secondary infection
    \item Lichenification and pigmentary changes from chronic scratching
    \item Disturbed sleep, fatigue, and irritability
    \item Anxiety and depression from persistent itching
    \item Reduced quality of life and work impairment
    \item Delayed diagnosis of underlying disease when itching is dismissed
\end{itemize}

\section*{Prognosis and Outcomes}
Most itching resolves with treatment of the cause and simple skin care. Dry skin responds to regular moisturising, eczema and hives respond to anti-inflammatory and antihistamine treatment, and infestations are cured with specific therapy. When an internal cause is found, treating it usually relieves the itch. Chronic and unexplained itching can be difficult, but dermatology review and modern treatments help most people achieve meaningful relief.

\section*{Prevention and Self-Care}
\begin{itemize}
    \item Moisturise daily, especially after bathing and in winter
    \item Use lukewarm rather than hot water for washing
    \item Avoid harsh soaps, perfumes, and fabric softeners that irritate
    \item Wear loose, soft, breathable fabrics
    \item Keep the home cool and use a humidifier in dry conditions
    \item Do not scratch: keep nails short and use cold compresses
    \item Avoid known triggers, including foods, medicines, and allergens
    \item See a doctor for persistent itching or itching without a rash
    \item Report itching in pregnancy promptly
\end{itemize}

\section*{Sources and Further Reading}
The following reputable sources provide further information for healthcare students and patients seeking reliable, current advice about itchy skin.
\begin{itemize}
    \item Healthdirect Australia. \textit{Itchy skin.} https://www.healthdirect.gov.au/
    \item Australasian College of Dermatologists. \textit{Itching and skin conditions.} https://www.dermcoll.edu.au/
    \item Eczema Association of Australasia. \textit{Eczema information and support.} https://www.eczema.org.au/
    \item Better Health Channel. \textit{Itching.} https://www.betterhealth.vic.gov.au/
    \item NHS. \textit{Itchy skin.} https://www.nhs.uk/conditions/itchy-skin/
    \item Mayo Clinic. \textit{Itchy skin.} https://www.mayoclinic.org/
\end{itemize}
